\uuid{g7FN}
\titre{Matrice d'une application linéaire/Matrice de passage}
\auteur{hanani ; mzari}
\organisation{exo7}
\licence{CC BY-NC-SA 4.0}
\datecreate{2020-01-20}
\niveau{L1}
\module{Algèbre}
\chapitre{Matrice}
\sousChapitre{Matrice et application linéaire}
\difficulte{1}
\theme{Matrice d'une application linéaire, Matrice de passage}
\exo7qcmid{lille-2.matrice-application.007}

\contenu{
\begin{question}
\qclassification{108.03}
\qtags{motcle=Matrice d'une application linéaire/Matrice de passage}
Dans $\Rr^3$, on considère la base canonique ${\cal {B}} =\{e_1,e_2,e_3\}$ et  la base ${\cal {B'}} =\{u_1,u_2,u_3\}$, où
$u_1=(1,1,-1), u_2=(0,2,1)$ et $u_3=(0,1,1)$. On notera $P$ la matrice de passage de la base ${\cal {B}}$ à la base ${\cal {B'}}$ et $Q$ la matrice de passage de la base ${\cal {B'}}$ à la base ${\cal {B}}$.
\vskip0mm
Définition : Soit $E$ un espace vectoriel muni de deux bases ${\cal {B}}$ et ${\cal {B}}'$.
La matrice de passage de la base ${\cal {B}}$ à la base  ${\cal {B}}'$ est la matrice de l'identité de $E$ de la base
${\cal {B}}'$ à la base  ${\cal {B}}$. Autrement dit, c'est la matrice dont la jième colonne
est constituée des coordonnées du jième vecteur de la base ${\cal {B}'}$ dans la base  ${\cal {B}}$.
\vskip2mm
Quelles sont les assertions vraies ?
\begin{answers}
\good{$P = \left(\begin{array}{rcc}
1&0&0\\
1&2&1\\
-1&1&1\\
\end{array}\right).$}
\bad{$Q = \left(\begin{array}{rcc}
1&0&0\\
1&2&1\\
-1&1&1\\
\end{array}\right).$}
\good{$Q = \left(\begin{array}{rcc}
1&0&0\\
-2&1&-1\\
3&-1&2\\
\end{array}\right).$}
\bad{$P$ est inversible et $P^{-1}=\left(\begin{array}{rcc}
1&0&0\\
1&2&1\\
-1&1&1\\
\end{array}\right).$}
\end{answers}
\begin{explanations} Une matrice de passage est inversible, puisque l'application linéaire associée est bijective. L'inverse de la matrice de passage de la base ${\cal {B}}$ à la base ${\cal {B}}'$ est la matrice de passage de la base ${\cal {B}}'$ à la base ${\cal {B}}$.\\
$P = \left(\begin{array}{rcc}
1&0&0\\
1&2&1\\
-1&1&1\\
\end{array}\right)$ et $
Q = P^{-1} = \left(\begin{array}{rcc}
1&0&0\\
-2&1&-1\\
3&-1&2\\
\end{array}\right)$.
\end{explanations}
\end{question}
}
